\section{Introduction}
\label{sec:introduction}

The Transformer architecture \citep{vaswani2017attention} has achieved rapid and widespread dominance in NLP. At its heart is a attention mechanism -- an inductive bias that connects each token in the input through a relevance weighted basis of every other token. Many papers have prodded and probed the Transformer, and in particular the attention sublayers, in an effort to better understand the architecture; see, for example, \citet{tenney-etal-2019-bert, vig-belinkov-2019-analyzing, clark-etal-2019-bert, voita-etal-2019-analyzing}. Although potentially limited in their effectiveness \citep{hewitt-liang-2019-designing}, these probes generally back the intuition that, by allowing higher order units to form out of compositions of the input, Transformer models can flexibly capture diverse syntactic and semantic relationships. 

In this work, we investigate whether simpler token mixing mechanisms can wholly replace the relatively complex self-attention layers in Transformer encoder architectures. We first replace the attention sublayer with two parameterized matrix multiplications -- one mixing the sequence dimension and one mixing the hidden dimension. Seeing promising results in this simple linear mixing scheme, we further investigate the efficacy of faster, structured linear transformations. Surprisingly, we find that the Fourier Transform, despite having no parameters at all, achieves nearly the same performance as dense linear mixing and scales very efficiently to long inputs, especially on GPUs (owing to the \O{N \log N} Fast Fourier Transform (FFT) algorithm). We call the resulting model FNet. 

While Fourier Transforms have previously been used to approximate or speed up computations in Convolutional Neural Networks~\citep{el2004fast, mathieu2014fast, highlander2016very, pratt2017fcnn, lin2018fft, chitsaz2020acceleration, goldberg2020rethinking}, Recurrent Neural Networks \citep{koplon1997using, zhang2000forenet, zhang2018learning},  Transformers~\citep{choromanski2020masked, tamkin2020language}, and MLP layers more generally~\citep{cheng2015exploration, moczulski2015acdc, sindhwani2015structured}, we believe our work is the first to wholly replace particular neural network sublayers with a Fourier Transform. This approach of viewing the Fourier Transform as a first class mixing mechanism is reminiscent of the MLP-Mixer~\citep{tolstikhin2021mlp} for vision, which replaces attention with MLPs; although in contrast to MLP-Mixer, FNet has no learnable parameters that mix along the spatial dimension.

Given the favorable asymptotic complexity of the FFT, our work also connects with the literature on ``long sequence'' or ``efficient'' Transformers, which aim to make the attention mechanism scale better via sparsity patterns~\citep{child2019generating, qiu2019blockwise, parmar2018image, beltagy2020longformer, ainslie2020etc, zaheer2020big, wang2020linformer, tay2020sparse, tay2020synthesizer, kitaev2020reformer, roy2021efficient, vyas2020fast, liu2018generating} or via linearization of the attention matrix~\citep{katharopoulos2020transformers, choromanski2020rethinking, peng2021random}. As we will show in our experiments, while some of those works achieve $O(N)$ scaling of attention, this complexity often hides large constants, which make them less scalable in practice than FNet.

The contributions of our paper are:
\begin{itemize}
    \item We show that simple linear transformations, including even (parameter-free) Fourier Transforms, along with standard MLPs in feed-forward layers, are competent at modeling diverse relationships in text. That such a simple linear transformation works at all is surprising, and suggests that, for at least some NLP problems, attention may not be the principal component driving the performance of Transformers.
    \item We introduce a new model, FNet, that uses the Fourier Transform as a mixing mechanism.  FNet offers an excellent compromise between speed, memory footprint, and accuracy, achieving $92\%$ and $97\%$, respectively, of the accuracy of BERT-Base and BERT-Large \citep{devlin2018bert} on the GLUE benchmark \citep{wang2018glue}, while training $80\%$ faster on GPUs and $70\%$ faster on TPUs. 
    \item We find that FNet hybrid models containing only two self-attention sublayers achieve $97-99\%$ of their BERT counterparts' accuracy on GLUE, while still running $40-70\%$ faster. This indicates that, while attention can improve accuracy, it may not be necessary to use in every layer.
    \item We demonstrate FNet scales very well to long inputs and offers a better compromise between speed and accuracy than the efficient Transformers evaluated on the Long-Range Arena (LRA) benchmark \citep{tay2020long}.
    Specifically, FNet achieves accuracy comparable to the most accurate efficient Transformer architectures but is significantly faster at both training and inference than all of the evaluated Transformer architectures across all sequence lengths on GPUs. On TPUs, FNet is faster for relatively shorter sequence lengths; for longer sequences, the only efficient Transformers that are faster than FNet on TPUs are less accurate on the LRA benchmark.
    Based on this, we argue that rather than seeking more efficient approximations of the attention, there may be more value in seeking out completely new mixing mechanisms.
\end{itemize}
